\documentclass{article}
\usepackage[utf8]{inputenc}

\usepackage{amsmath, amssymb, amsthm, amsfonts}
\usepackage{subcaption, enumerate}
\usepackage{bbold}
\usepackage{polynom}
\usepackage{xfrac}
\usepackage{graphicx}
\usepackage{mathtools}
\usepackage{enumitem}

\usepackage{tikz}

\newtheorem{theorem}{Theorem}
\newtheorem{corot}{Corollary}[theorem]
\newtheorem{lemma}{Lemma}
\newtheorem{corol}{Corollary}[lemma]
\newtheorem{proposition}{Proposition}
\newtheorem{corop}{Corollary}[proposition]

\theoremstyle{definition}
\newtheorem{definition}{Definition}
\newtheorem{exampled}{Example}[definition]
\newtheorem{examples}{Example}[section]
\newtheorem*{notation}{Notation}
\newtheorem*{exercise}{Exercise}

\newcommand{\dab}
{
	\begin{tikzpicture}[scale=0.04]
		% head
		\draw[thick] (0, 0) circle (3cm);

		% bent arm
		\draw[thick, rotate around={15 : (0, 0)}] (-3, 0) -- (-3.7, -2) -- (4.1, -1.1);

		% extended arm
		\draw[thick, rotate=20] (3, 0) -- (7, 2);
	\end{tikzpicture}
}
\newcommand{\st}{\text{ s.t. }}

\makeatletter
\renewcommand*\env@matrix[1][*\c@MaxMatrixCols c]{%
  \hskip -\arraycolsep
  \let\@ifnextchar\new@ifnextchar
  \array{#1}}
\makeatother

\usepackage{fullpage}

\title{Topology \\ with Dr Nevard}
\author{Laura Gallo}
\date{November 16, 2020}

\begin{document}
\maketitle

\section{Interiors and Exteriors}
\begin{definition}
	Given $S \subset X$, the \emph{interior} of $S$ is:
	$$S^\circ=\bigcup_\text{$U \subset S$, $U$ open}U$$
\end{definition}
\begin{exampled}
	For $S=[0, 1]$ in the standard topology for $\mathbb{R}$, the interior of $S$ is $S^\circ=(0, 1)$.
\end{exampled}
\begin{exampled}
	For $\mathbb{Q}$ in the standard topology for $\mathbb{R}$, $Q^\circ=\emptyset$.
\end{exampled}

\begin{definition}
	The \emph{exterior} of a set $S$ is $\text{ext } S =(S^c)^\circ$.
\end{definition}
\begin{exampled}
	$\text{ext } [0, 1] = (-\infty, 0) \cup (1, \infty)$.
\end{exampled}

\section{Closed Sets}
\begin{definition}
	A \emph{closed set} is a set whose complement is an open set; ie, a set $C$ is closed if and only if $C^c$ is open. \\
	Note: a set can be both open and closed, and also neither open nor closed (example: $\emptyset$ and $X$ for a topology on $X$ are both open and closed)
\end{definition}
\noindent
We have $\bigcap_\alpha U_\alpha^c=(\bigcup_\alpha U_\alpha)^c$, so closed sets are closed under intersections because open sets are closed under unions. Note, however, that open sets are not necessarily closed under intersections, and vice versa. For example, take $S_n=[-\frac{1}{n}, \frac{n+1}{n}]$ in the standard topology on $\mathbb{R}$:
\begin{align*}
	\bigcap_{n=1}^\infty S_n = [0, 1] && \text{(closed $\rightarrow$ closed)} \\
	\bigcup_{n=1}^\infty S_n = [0, 1] && \text{(open $\rightarrow$ closed)} \\
\end{align*}

\section{Closure}
\begin{definition}
	The \emph{closure} of $S$ is:
	$$\overline{S}=\bigcap_\text{$C \supset S$, $C$ closed} C$$
\end{definition}
\begin{exampled}
	We claim $\overline{\mathbb{Q}} = \mathbb{R}$.

	\begin{proof}
		Assume $\overline{\mathbb{Q}} \neq \mathbb{R}$; ie, there is some $x \in \mathbb{R}$ such that $x \notin \overline{\mathbb{Q}}$. \\
		We can find a closed set $C \supset \mathbb{Q}$ such that $x \notin C$. \\
		Let $U=C^c$. $U$ is an open set by definition. \\
		We can say $(x-\delta, x+\delta) \subset U$. \\
		This is a contradiction since this interval contains a rational number. \\
		Therefore, $\overline{\mathbb{Q}} = \mathbb{R}$. \dab
	\end{proof}
\end{exampled}

\begin{definition}
	If $S \subset X$ such that $\overline{S}=X$, then $S$ is \emph{dense} in $X$.
\end{definition}

\section{Limit Points and Isolated Points}
\begin{definition}
	A \emph{limit point} of a set $S$ is an element of $S$, $x$, such that every open set containing $x$ contains other points of $S$
\end{definition}
\begin{exampled}
	$S=\{\frac{1}{n} : n \in\mathbb{N}\}$, 0 is a limit point.
\end{exampled}

\begin{definition}
	An \emph{isolated point} of a set $S$ is a point that is not a limit point.
\end{definition}

\begin{proposition}
	A set is closed if and only if it contains all of its limit points
\end{proposition}

\section{Boundaries}
\begin{definition}
	The \emph{boundary} of a set $S$ is $\partial S = \overline{S} \cap \overline{S^c}$.
\end{definition}
\begin{exampled}
	$\partial \mathbb{Q}=\mathbb{R}$
\end{exampled}

\begin{exercise}
	Show that the boundary of $S=\{(x,y)\in\mathbb{R}^2 : x^2+y^2<1\}$ is $\{(x, y)\in\mathbb{R}^2 : x^2 + y^2 = 1\}$.
\end{exercise}

\section{Generalisation of the Definition of Limits}
We want to be able to say things such as $\lim_{x\rightarrow a} f(x)=L$ within a topological space. From least abstract to most abstract, we can state the following definitions:
\begin{itemize}
	\item We can make $f(x)$ arbitrarily close to $L$ by making $x$ sufficiently close to $a$, but not necessarily equal to $a$.
	\item In normed spaces $X$, $Y$:
		$$\forall \varepsilon > 0 \exists \delta > 0 \st \forall x \in X, 0<||x-a||<\delta \Rightarrow ||f(x)-L||<\varepsilon$$
	\item In metric spaces $X$, $Y$:
		$$\forall \varepsilon > 0 \exists \delta > 0 \st \forall x \in X, 0<d(x, a)<\delta \Rightarrow d(f(x), L)<\varepsilon$$
	\item In topological spaces $X$, $Y$:
		$$\forall V \subset Y \text{ open containing $L$} \exists U \subset X \text{ open containing x} \st z\in U, x \neq z \Rightarrow f(z) \in V$$
\end{itemize}

\end{document}

